\documentclass[a4paper,12pt]{article}
\usepackage[T2A]{fontenc}
\usepackage[utf8]{inputenc}
\usepackage[russian]{babel}
\usepackage{amsmath,amssymb,mathtools}
\renewcommand{\familydefault}{\sfdefault}
\usepackage{icomma}
\usepackage{geometry}
\geometry{left=19mm,right=19mm,top=19mm,bottom=19mm}
\usepackage{microtype}
\usepackage{xcolor}
\usepackage{tikz}
\usetikzlibrary{arrows.meta,positioning,shapes.geometric,calc}
\usepackage{tabularx,booktabs}
\usepackage{enumitem}
\usepackage{fancyhdr}
\usepackage{setspace}
\usepackage{hyperref}
\onehalfspacing

\definecolor{accent}{HTML}{1B7F79}
\definecolor{darkaccent}{HTML}{135C58}
\definecolor{textgray}{HTML}{343B3D}
\definecolor{softgray}{HTML}{F3F6F5}
\definecolor{warn}{HTML}{915A35}
\hypersetup{colorlinks=true,linkcolor=darkaccent,urlcolor=darkaccent}
\setlength{\parindent}{0pt}
\setlist[itemize]{leftmargin=5mm,itemsep=0.6mm,topsep=0.7mm}
\setlist[enumerate]{leftmargin=7mm,itemsep=2mm,topsep=1mm}
\setlength{\headheight}{15pt}
\pagestyle{fancy}
\fancyhf{}
\fancyhead[L]{\small\color{textgray}Тригонометрия и повторение ЕГЭ}
\fancyhead[R]{\small\color{textgray}Николь Саркисьянц}
\fancyfoot[C]{\small\color{textgray}\thepage}
\renewcommand{\headrulewidth}{0.3pt}
\newcommand{\key}[1]{\textcolor{darkaccent}{\textbf{#1}}}
\newcommand{\attention}[1]{\textcolor{warn}{\textbf{#1}}}
\newcommand{\R}{\mathbb{R}}

\tikzset{
 procbox/.style={draw=accent,rounded corners=2mm,thick,text width=108mm,align=center,minimum height=13mm,fill=white,inner sep=3mm},
 choice/.style={draw=accent,diamond,aspect=3.2,text width=68mm,align=center,fill=softgray,inner sep=1mm,thick},
 terminus/.style={draw=accent,ellipse,align=center,minimum width=58mm,minimum height=10mm,fill=softgray,thick},
 linearrow/.style={-{Latex[length=2.7mm]},draw=darkaccent,thick}
}

\begin{document}
\thispagestyle{empty}
\begin{center}
\vspace*{19mm}
{\Large\color{darkaccent}\textbf{Чек-лист}}\\[6mm]
{\huge\bfseries Тригонометрические уравнения}\\[3mm]
{\Large Формулы, однородность и ограничения}\\[13mm]
{\large\bfseries Николь Саркисьянц}\\[3mm]
{\normalsize Профильная математика, ЕГЭ}\\[14mm]
\begin{minipage}{0.81\linewidth}
\centering\color{textgray}
Повторение ключевых приёмов: переход к половинному углу,
решение однородных уравнений, точная работа с ограниченностью
функций, логарифмическими выражениями и квадратной моделью движения.
\end{minipage}
\vfill
{\large 14.10.2026}\\[8mm]
{\small Лёвин Артём Александрович эксклюзивно для Николь Саркисьянц}
\vspace{10mm}
\begin{tikzpicture}[x=1cm,y=1cm]
\draw[accent,line width=1.1pt] (-6,0) .. controls (-4,0.9) and (-2,-0.8) .. (0,0)
 .. controls (2,0.8) and (4,-0.7) .. (6,0.16);
\end{tikzpicture}
\end{center}
\clearpage

\section*{1. Узнаём вид тригонометрического уравнения}
\begin{itemize}
\item $\square$ Вижу $a\sin x+b\cos x=c$ и знаю, что его удобно преобразовать через формулы двойного угла для угла $x/2$.
\item $\square$ Отличаю однородное уравнение первой степени $a\sin u+b\cos u=0$ от второй степени $a\sin^2u+b\sin u\cos u+c\cos^2u=0$.
\item $\square$ Перед делением на $\cos u$ или $\cos^2u$ отдельно проверяю случай $\cos u=0$.
\item $\square$ После замены $t=\tg u$ возвращаюсь к углу: $u=\arctg t+\pi k$, $k\in\mathbb Z$.
\end{itemize}

\subsection*{Формулы, необходимые в ходе решения}
\[
\begin{aligned}
\sin 2u&=2\sin u\cos u,\\
\cos 2u&=\cos^2u-\sin^2u,\\
1&=\sin^2u+\cos^2u.
\end{aligned}
\]
\key{Главная идея:} в уравнении с углом $x$ положите $u=x/2$; в уравнении с углом $2x$ положите $u=x$.

\subsection*{Пример 1. $\sin x+\cos x=1$}
Полагаем $u=x/2$ и преобразуем каждое слагаемое:
\[
2\sin u\cos u+\cos^2u-\sin^2u
=\sin^2u+\cos^2u.
\]
Сокращая одинаковые слагаемые, получаем
\[
2\sin u(\cos u-\sin u)=0.
\]
Поэтому $\sin u=0$ или $\cos u-\sin u=0$.
Во второй ветви $\cos u\ne0$, поэтому $\tg u=1$.
\[
\begin{aligned}
&u=\pi k &&\Longrightarrow\quad x=2\pi k,\\
&u=\frac\pi4+\pi k &&\Longrightarrow\quad x=\frac\pi2+2\pi k.
\end{aligned}\qquad k\in\mathbb Z
\]

\attention{Важно:} не делите всё исходное уравнение на $\cos u$: при $\cos u=0$ могли бы потеряться допустимые решения других ветвей.

\subsection*{Другой путь: формулы приведения и суммы}
\[
\cos x=\sin\!\left(\frac\pi2-x\right),\qquad
\sin\alpha+\sin\beta
=2\sin\frac{\alpha+\beta}{2}\cos\frac{\alpha-\beta}{2}.
\]
Отсюда $\sin x+\cos x=\sqrt2\cos(x-\pi/4)$, и мы получаем то же множество корней.
\clearpage

\section*{2. Однородное уравнение второй степени}
\subsection*{Пример 2. $3\sin x+4\cos x=5$}
Переходим к $u=x/2$:
\[
6\sin u\cos u+4(\cos^2u-\sin^2u)
-5(\sin^2u+\cos^2u)=0.
\]
После приведения подобных:
\[
6\sin u\cos u-\cos^2u-9\sin^2u=0.
\]
При $\cos u=0$ левая часть равна $-9$, поэтому такие углы не являются решениями. Делим на $\cos^2u$:
\[
6\tg u-1-9\tg^2u=0
\quad\Longleftrightarrow\quad (3\tg u-1)^2=0.
\]
Следовательно, $\tg(x/2)=1/3$ и
\[
\boxed{x=2\arctg\frac13+2\pi k,\quad k\in\mathbb Z.}
\]
\key{Самопроверка по ограниченности:} $3\sin x+4\cos x\le5$; равенство возможно при $\sin x=3/5$, $\cos x=4/5$.

\subsection*{Пример 3. $24\sin 2x-7\cos 2x=15$}
Разложим $\sin 2x$, $\cos 2x$ и $15=15(\sin^2x+\cos^2x)$:
\[
48\sin x\cos x-22\cos^2x-8\sin^2x=0,
\]
\[
24\sin x\cos x-11\cos^2x-4\sin^2x=0.
\]
При $\cos x=0$ исходное уравнение даёт $7=15$, что невозможно.
Для $\cos x\ne0$ делим на $\cos^2x$:
\[
4t^2-24t+11=0,\qquad t=\tg x.
\]
\[
t_1=\frac12,\qquad t_2=\frac{11}{2};\qquad
\boxed{x=\arctg\frac12+\pi k\;\;\text{или}\;\;x=\arctg\frac{11}{2}+\pi k.}
\]

\subsection*{Ещё две формулы преобразования сумм}
\[
\begin{aligned}
\cos\alpha+\cos\beta
 &=2\cos\frac{\alpha+\beta}{2}\cos\frac{\alpha-\beta}{2},\\
\cos\alpha-\cos\beta
 &=-2\sin\frac{\alpha+\beta}{2}\sin\frac{\alpha-\beta}{2}.
\end{aligned}
\]
\key{Признак:} вид $a\sin^2u+b\sin u\cos u+c\cos^2u=0$ подсказывает замену $t=\tg u$ после проверки делителя.
\clearpage

\section*{3. Ограниченность: сначала оцениваем выражение}
Для любого $x$ верно $-1\le\sin x\le1$ и $-1\le\cos x\le1$.
В частности, $0\le\sin^4x\le1$, но \attention{нельзя} утверждать, что $(\sin x-2)^4\le1$.

\subsection*{Пример 4. $\sin^4x+(\sin x-2)^4=2$}
Положим $u=\sin x-1$. Тогда $\sin x=u+1$, а $\sin x-2=u-1$. Получаем
\[
(u+1)^4+(u-1)^4=2
\ \Longleftrightarrow\
2u^4+12u^2+2=2.
\]
Следовательно, $2u^2(u^2+6)=0$, откуда $u=0$, $\sin x=1$:
\[
\boxed{x=\frac\pi2+2\pi k,\quad k\in\mathbb Z.}
\]
\key{Смысл приёма:} сумма равна своему наименьшему значению $2$ только при $u=0$. Проверяйте границы \emph{всех} слагаемых прежде, чем приравнивать каждое к максимуму.

\section*{4. Логарифм и тригонометрия в одной задаче}
\begin{itemize}
\item $\square$ Учитываю условия $a>0$, $a\ne1$, аргумент логарифма $>0$.
\item $\square$ Помню $\log_a(MN)=\log_aM+\log_aN$, $\log_a M^r=r\log_aM$ при допустимых аргументах.
\item $\square$ Использую $\log_a b=1/\log_b a$ и $\log_a a=1$.
\item $\square$ В уравнении $\log_a F(x)=c$ перехожу к $F(x)=a^c$; правая часть положительна, поэтому требование $F(x)>0$ выполняется автоматически для найденных корней.
\end{itemize}

\subsection*{Пример 5. Степени и смена основания}
Если $\log_b a=-2$, то
\[
\begin{aligned}
\log_a(b^6\sqrt a)
&=6\log_a b+\tfrac12\log_a a
=\frac6{\log_b a}+\frac12
=\boxed{-\frac52}.
\end{aligned}
\]
Здесь предполагается $a,b>0$, $a,b\ne1$.

\clearpage
\subsection*{Пример 6. Логарифм от тригонометрического выражения}
\[
\log_2\!\left(\sqrt2\cos\!\left(x-\frac\pi2\right)
-\cos\!\left(x-\frac\pi4\right)\right)=-\frac12.
\]
По определению логарифма выражение в скобках равно $2^{-1/2}=\sqrt2/2>0$.
Используя формулы разности углов,
\[
\cos\!\left(x-\frac\pi2\right)=\sin x,\qquad
\cos\!\left(x-\frac\pi4\right)=\frac{\sqrt2}{2}(\cos x+\sin x),
\]
получаем $\sin x-\cos x=1$. Следовательно,
\[
\boxed{x=\frac\pi2+2\pi k\;\;\text{или}\;\;x=\pi+2\pi k,\quad k\in\mathbb Z.}
\]
\vspace{4mm}
\section*{5. Квадратичная модель движения}
Высота мяча (в метрах) через $t$ секунд после броска:
\[
h(t)=-10t^2+16t+1,7,\qquad t\ge0.
\]
\key{Вопрос:} сколько секунд мяч находился на высоте \emph{не менее} $4,5$~м?
Решаем неравенство, учитывая физический смысл времени:
\[
-10t^2+16t+1,7\ge4,5
\ \Longleftrightarrow\
25t^2-40t+7\le0.
\]
Корни соответствующего уравнения:
\[
t_1=\frac15=0,2,\qquad t_2=\frac75=1,4.
\]
Парабола $25t^2-40t+7$ направлена ветвями вверх, поэтому
\[
t\in[0,2;\,1,4],\qquad \boxed{\Delta t=1,4-0,2=1,2\ \text{с}.}
\]
\key{Важное различие:} $t_1$ и $t_2$ --- моменты пересечения заданной высоты на подъёме и при падении. Ответом на вопрос «как долго?» является $t_2-t_1$.

\subsection*{Перед записью ответа}
Проверьте \key{метод}, \key{углы} в тождествах, \key{нулевой делитель},
\key{ОДЗ}, \key{подстановку} корней и \key{формат ответа}
($k\in\mathbb Z$, единицы измерения, требуемая величина).
\clearpage

\thispagestyle{plain}
\begin{center}{\Large\color{darkaccent}\bfseries Алгоритм решения тригонометрического уравнения}\end{center}
\vspace{6mm}
\begin{center}
\begin{tikzpicture}[node distance=10mm]
\node (start) [terminus] {Уравнение $a\sin u+b\cos u=c$};
\node (convert) [procbox,below=13mm of start] {Представьте $\sin u$ и $\cos u$ через угол $v=u/2$:\\
$\sin u=2\sin v\cos v$, $\cos u=\cos^2v-\sin^2v$};
\node (hom) [procbox,below=of convert] {Замените число $c$ выражением $c(\sin^2v+\cos^2v)$; приведите подобные слагаемые};
\node (check) [procbox,below=of hom] {Получите однородное уравнение.\\ Отдельно проверьте $\cos v=0$ и сохраните подходящие корни};
\node (divide) [procbox,below=of check] {При $\cos v\ne0$ разделите на $\cos^2v$ и положите $t=\tg v$};
\node (solve) [procbox,below=of divide] {Решите квадратное уравнение по $t$, затем $v=\arctg t+\pi k$};
\node (finish) [terminus,below=of solve] {Вернитесь к $x$; объедините и проверьте корни};
\draw[linearrow] (start)--(convert);
\draw[linearrow] (convert)--(hom);
\draw[linearrow] (hom)--(check);
\draw[linearrow] (check)--(divide);
\draw[linearrow] (divide)--(solve);
\draw[linearrow] (solve)--(finish);
\end{tikzpicture}
\end{center}
\clearpage

\section*{Тренировочный блок --- Путь А}
\textcolor{textgray}{Решите 10 заданий. Оформляйте ключевые преобразования, ограничения и общий вид корней. Решения в этот блок не включены.}
\begin{enumerate}[label=\textbf{\arabic*.},itemsep=6mm]
\item Решите уравнение $\sin x+\cos x=1$.
\item Решите уравнение $\sin x-\cos x=1$.
\item Решите уравнение $\sin 2x-2\cos 2x=1$.
\item Решите уравнение $3\sin x+4\cos x=5$.
\item Решите уравнение $24\sin 2x-7\cos 2x=15$.
\item Известно, что $\log_b a=-2$, где $a,b>0$, $a,b\ne1$. Вычислите $\log_a(b^6\sqrt a)$.
\item Высота мяча задаётся формулой $h(t)=-10t^2+16t+1,7$ (в метрах, $t$ в секундах). Сколько секунд мяч находится на высоте не менее $4,5$~м?
\item Решите уравнение
\[
\log_2\!\left(\sqrt2\cos\!\left(x-\frac\pi2\right)-\cos\!\left(x-\frac\pi4\right)\right)=-\frac12.
\]
\item Решите уравнение
\[
\log_2\!\left(\sqrt2\cos\!\left(x-\frac\pi2\right)+\cos\!\left(x+\frac\pi4\right)\right)=-\frac12.
\]
\item Решите уравнение $\sin^4x+(\sin x-2)^4=2$; обоснуйте, почему других корней нет.
\end{enumerate}
\clearpage

\section*{Ответы к тренировочному блоку}
\textcolor{textgray}{Во всех тригонометрических ответах $k\in\mathbb Z$.}\par\medskip\noindent
\renewcommand{\arraystretch}{1.55}
\begingroup\small
\begin{tabularx}{\linewidth}{@{}p{9mm}X@{}}
\toprule
\textbf{№}&\textbf{Ответ}\\
\midrule
1 & $x=2\pi k$ или $x=\dfrac\pi2+2\pi k$.\\
2 & $x=\dfrac\pi2+2\pi k$ или $x=\pi+2\pi k$.\\
3 & $x=\dfrac\pi4+\pi k$ или $x=-\arctg3+\pi k$.\\
4 & $x=2\arctg\dfrac13+2\pi k$.\\
5 & $x=\arctg\dfrac12+\pi k$ или $x=\arctg\dfrac{11}{2}+\pi k$.\\
6 & $-\dfrac52$.\\
7 & $1,2$~с.\\
8 & $x=\dfrac\pi2+2\pi k$ или $x=\pi+2\pi k$.\\
9 & $x=2\pi k$ или $x=\dfrac\pi2+2\pi k$.\\
10 & $x=\dfrac\pi2+2\pi k$.\\
\bottomrule
\end{tabularx}\par
\endgroup
\vspace{7mm}\noindent
\textcolor{textgray}{При сверке проверьте не только окончательные корни, но и промежуточные ограничения: в заданиях 3--5 --- деление на косинус, в заданиях 8--9 --- положительность аргумента логарифма.}
\end{document}